\section*{第 \ref{ch:b3:harmonic-oscillator} 章——量子谐振子}

\begin{solution}{exo:b3:harmonic-oscillator:1}
(a) 展开后，$\hat x^2$ 与 $\hat p^2$ 项在
$\hat a\hat a^\dagger$ 与 $\hat a^\dagger\hat a$ 之间相消，剩下
$(m\omega/2\hbar)(-2\iu/m\omega)[\hat x, \hat p] =
(-\iu/\hbar)(\iu\hbar) = 1$。(b) $\hat a^\dagger\hat a = (m\omega\hat x^2/2\hbar) +
(\hat p^2/2m\hbar\omega) - \tfrac12$：乘以 $\hbar\omega$。(c)
$\hat x = \sqrt{\hbar/2m\omega}\,(\hat a + \hat a^\dagger)$，$\hat p
= \iu\sqrt{m\hbar\omega/2}\,(\hat a^\dagger - \hat a)$。(d) $\hat
N^\dagger = \hat a^\dagger\hat a = \hat N$；$\hat a$ 不是
厄米的，其“\omterm{def:b3:quantum-formalism:observable}{本征值}” $\alpha$ 为复数——没有
测量装置能测得它们。
\end{solution}

\begin{solution}{exo:b3:harmonic-oscillator:2}
(a) $\hat x$ 与 $\hat p$ 使 $n$ 改变 $\pm1$：对角元
为零。(b) $\langle\hat x^2\rangle = (\hbar/2m\omega)(2n + 1)$，
$\langle\hat p^2\rangle = (m\hbar\omega/2)(2n + 1)$（$\hat
a\hat a^\dagger + \hat a^\dagger\hat a$ 项）。(c) $\Delta x\Delta p
= (n + \tfrac12)\hbar$：仅基态取最小值。(d) 两者
平均值均为 $E_n/2$：经典
振子的能量均分定理，逐能级成立。
\end{solution}

\begin{solution}{exo:b3:harmonic-oscillator:3}
(a) $\Delta x = \qty{3.4}{pm}$：键长的百分之三——一个
分子即便在零温下也是略显模糊的物体。(b)
$\qty{14}{pm}$，超过十分之一埃：氢是
任何晶体中最模糊的原子，中子散射可直接
观测到它。(c) $\qty{4.6e-20}{m}$——比镜子的
尺度低二十五个数量级，但 LIGO 的读出仍能达到它。(d) 相对而言：
分子 $3\%$，氢 $\sim15\%$，镜子 $\sim10^{-19}$：
“量子”并非关乎尺寸小，而是关乎 $\hbar\omega$
相对于其他一切——然而只要足够精细，即使四十
千克的物体也会显露它们的底噪。
\end{solution}

\begin{solution}{exo:b3:harmonic-oscillator:4}
(a) $\bra m\hat x\ket n = \sqrt{\hbar/2m\omega}\,(\sqrt n\,
\delta_{m,n-1} + \sqrt{n + 1}\,\delta_{m,n+1})$。(b) 光与
$\propto\hat x$ 的相互作用只能跃迁一个能级：$\Delta n =
\pm1$，每个模式对应一个红外频率。(c) CO 的振动调制
一个非零的偶极矩；N$_2$ 的对称电荷云在
任何伸缩下都不产生偶极矩：红外非活性。(d) 泛音的存在，只因
真实势能并非严格二次：非简谐性使能级相互混合，
并微弱地允许 $\Delta n = 2$——并使高能级
彼此靠得更近，正如真实光谱所示。
\end{solution}

\begin{solution}{exo:b3:harmonic-oscillator:5}
(a) $\varphi_0' = -(m\omega/\hbar)x\varphi_0$：正文中的归一化
高斯函数。(b) $\varphi_1 \propto x\eu^{-m\omega
x^2/2\hbar}$；$\varphi_2 \propto (2m\omega x^2/\hbar -
1)\eu^{-m\omega x^2/2\hbar}$。(c) $\varphi_0$ 偶，$\varphi_1$ 奇：
重叠积分的被积函数为奇。(d) 零、一、两个节点——振荡定理的
缩影。
\end{solution}

\begin{solution}{exo:b3:harmonic-oscillator:6}
(a) 令 $x = \hbar\omega/k_{\text{B}}T$：$\langle n\rangle =
\sum n\eu^{-nx}/\sum\eu^{-nx} = 1/(\eu^x - 1)$。(b) CO 在
\qty{300}{K}：$x = 10.4$，$\langle n\rangle \approx \num{3e-5}$——
冻结；在 \qty{2000}{K}：$x = 1.57$，$\langle n\rangle = 0.26$——
苏醒。(c) $\langle E\rangle = \hbar\omega(\langle n\rangle +
\tfrac12) \to k_{\text{B}}T$，当 $x \ll 1$ 时；$\to \tfrac12\hbar
\omega$，当 $x \gg 1$ 时。(d) $\eu^x = 11$：$x = 2.4$，$T =
\hbar\omega/2.4k_{\text{B}} \approx \qty{400}{K}$——地球的
大气层使二氧化碳的弯曲振动模式部分保持激发。
\end{solution}

\begin{solution}{exo:b3:harmonic-oscillator:7}
(a) $\mathcal P(n) = |c_n|^2 = \eu^{-|\alpha|^2}|\alpha|^{2n}/n!$：
泊松分布。(b) 泊松分布的均值与方差均为 $|\alpha|^2$。
(c) 每秒 $\num{3.2e15}$ 个光子；$\Delta N/\langle N\rangle =
1/\sqrt{\langle N\rangle} \approx \num{1.8e-8}$。(d) 光电流
继承泊松分布的涨落：噪声与信号之比
$\propto 1/\sqrt N$——散粒噪声底，在微弱
光下清晰可辨，强光下可忽略。
\end{solution}

\begin{solution}{exo:b3:harmonic-oscillator:8}
(a) 每一项乘上 $\eu^{-\iu n\omega t}$（略去整体相位）：
其和仍是对应于 $\alpha\eu^{-\iu\omega t}$ 的\omterm{def:b3:harmonic-oscillator:coherent}{相干态}。(b)
$\langle\hat x\rangle = \sqrt{2\hbar/m\omega}\,
\operatorname{Re}\,\alpha(t)$：频率为 $\omega$ 的纯余弦，振幅
和相位由 $\alpha$ 决定——与经典完全相同。(c) 本征态的
密度从不移动；我们所见的摆动，是
许多能级的相干叠加，其相位协同使之摆动。(d) $E
\approx \qty{2e-3}{J}$，$\hbar\omega \approx \qty{6.6e-34}{J}$：$n
\sim \num{3e30}$，$|\alpha| \sim \num{2e15}$——宏观运动
即是具有天文量级量子数的相干性。
\end{solution}

\begin{solution}{exo:b3:harmonic-oscillator:9}
(a) $\omega = 1/\sqrt{LC} = \qty{1.6e10}{rad/s}$（$f =
\qty{2.5}{GHz}$）；$\hbar\omega = \qty{10}{\micro eV}$。(b)
$\hbar\omega/k_{\text{B}} = \qty{0.12}{K}$：远低于约
\qty{100}{mK}。(c) 在 \qty{20}{mK} 下，热激发
$\eu^{-6} \sim \num{2e-3}$：电路处于 $\ket0$，可用作
量子比特——室温会将量子淹没于 $10^3$
个热量子之中。(d) $\Delta q = \sqrt{\hbar\omega C/2} =
\qty{5.8e-19}{C}$（约几个电子电荷）；$\Delta V = \Delta q/C
\approx \qty{1.4}{\micro V}$ 的不可消除的本底噪声。
\end{solution}

\begin{solution}{exo:b3:harmonic-oscillator:10}
(a) 在简正坐标下，\omterm{def:b3:hamiltonian-mechanics:hamiltonian}{哈密顿量}分裂为两个
谐振子，满足 $m\omega_\pm^2 = m\omega_0^2 \pm \lambda$。(b)
$\sqrt{1 + u} + \sqrt{1 - u} \approx 2 - u^2/4$：$\Delta E =
\tfrac\hbar2(\omega_+ + \omega_- - 2\omega_0) =
-\hbar\omega_0\lambda^2/8m^2\omega_0^4 =
-\hbar\lambda^2/8m^2\omega_0^3$，且 $\lambda \propto 1/R^3$ 给出
$1/R^6$。(c) 它不需要永久电偶极子——只需
各原子电荷云的零点\emph{涨落}，耦合使其
相关，以致吸引压过排斥。(d) 涨落
量是基态的瞬时偶极矩；
若 $\hbar = 0$，基态将静止
且无偶极矩，范德瓦耳斯“胶水”——壁虎、液化
气体与大部分软物质——都将消失。
\end{solution}

\begin{solution}{exo:b3:harmonic-oscillator:11}
(a) $hf = \qty{4.1}{neV}$：光学谱线必须窄于
一兆赫兹——只有寿命长的“时钟”跃迁才符合要求。(b)
红边带光子使离子的电子能级升高，同时
\emph{移除}一个运动量子；随后的衰变
平均地以载波频率返还电子能量：每
个循环提取 $\hbar\omega_{\text{囚禁}}$ 的运动能量。(c) 若从 $n = 0$ 出发，
则已无可移除之物：红边带的消失
是基态的一个无背景干扰的判据。(d)
$\langle n\rangle = 0.05$: $T = \hbar\omega/k_{\text{B}}\ln21
\approx \qty{16}{\micro K}$.
\end{solution}

\begin{solution}{exo:b3:harmonic-oscillator:12}
(a) $\dd t = \dd x/|v|$，其中 $v = \omega\sqrt{A^2 - x^2}$，积分区间为
半周期 $T/2 = \pi/\omega$。(b) 二者在
经典允许区域内的局部平均上一致；在转折点
处必然不同（经典发散，量子为有限峰），
越过转折点后亦不同（禁戒区域内的量子拖尾）。(c) $\Delta E/E = 1/(n +
\tfrac12)$。(d) $n \sim \num{3e30}$：间距为 $10^{30}$ 分之一
——没有任何可想象的测量能分辨该能级阶梯，力学
看起来是连续的。
\end{solution}

\begin{solution}{pb:b3:harmonic-oscillator:1}
\textbf{1.} 三个平动；\emph{两个}转动（绕
分子轴的自旋不移动任何原子）；$9 - 5 = 4$ 个振动。
\textbf{2.} 沿轴的伸缩模式（对称：两个 O 同时
外移；不对称：C 在两者之间往返）；弯曲可以发生在
两个相互垂直的平面中的任一个——频率相同，
二重简并。
\textbf{3.} 不对称伸缩与弯曲使电荷
中心相互分开：振荡偶极子，具有红外活性。对称
伸缩始终使分子的偶极矩为零：
无红外活性。
\textbf{4.} $15.0$、$4.26$ 与 \qty{7.2}{\micro m}：均属红外；
\qty{15}{\micro m} 深处于对应地球温度的\emph{热}红外
区段。
\textbf{5.} 同核分子在振动时永不产生偶极矩；
大气主体对红外透明，整个温室效应
只依赖于痕量多原子气体。
\textbf{6.} 水的吸收带之间留有 $8$--$13\,\unit{\micro m}$
窗口；二氧化碳的 \qty{15}{\micro m} 吸收带正处于其
边缘，水在此竞争较弱——这种气体守护着一扇
水虚掩着的门。
\textbf{7.} $\hbar\omega_2 = 667 \times \num{1.24e-4} =
\qty{82.7}{meV}$; $E_n = (n + \tfrac12) \times \qty{82.7}{meV}$.
\textbf{8.} $\eu^{-82.7/24.8} = \eu^{-3.33} \approx 3.6\%$ 位于 $n =
1$；$0.13\%$ 位于 $n = 2$：阶梯处于轻微而持久的点亮之中。
\textbf{9.} $\lambda = hc/\hbar\omega \to \qty{15.0}{\micro m}$；
相同的间距既决定吸收也决定发射——往返
都是同一波长。
\textbf{10.} 间距相等：每一步 $n \to n + 1$ 所需的
光子相同，因此一条谱线即可覆盖整个热布居。
\textbf{11.} 每条振动谱线分裂为许多条
转动-振动谱线，分布于 $\sim\qty{1}{\micro m}$ 范围内：
一旦带中心已完全不透明，只有这一\emph{宽度}
才能提供新的吸收——吸收带的翼部正是额外气体仍在起作用之处。
\textbf{12.} 碰撞的占优程度高出九个数量级：
光子能量在纳秒内即与 N$_2$ 和 O$_2$ 分享
——被吸收的辐射转化为\emph{空气的热量}。
\textbf{13.} 同样的碰撞反向进行时，空气不断
向 $n = 1$ 输送分子，这些分子以 \qty{15}{\micro m} 向各方向
辐射——包括\emph{向下}：天空本身
在地面处发出红外辉光。
\textbf{14.} $2898/5800 = \qty{0.50}{\micro m}$：日光峰值位于
可见光区，远离 \qty{15}{\micro m}——二氧化碳让
太阳光得以进入。
\textbf{15.} $2898/288 = \qty{10.1}{\micro m}$：地球
辐射峰值位于十微米，而 \qty{15}{\micro m} 位于其宽阔的
肩部，承载了出射功率的相当大部分。
\textbf{16.} 日光大多无阻地进入并加热
地面。地面以热红外形式再辐射。在
\qty{15}{\micro m} 处，该辐射在数米内即被吸收并
热化。受热空气向上和向下同时再辐射，其中
向下的一半是地表额外的收入：地表必须升温，
直到支出与收入相等。
\textbf{17.} 发射随温度急剧增长：
从 \qty{220}{K} 高度逸出的辐射，其功率
远小于地表直接发出的功率——该吸收带
在行星的出射光谱中呈较暗的区域。
\textbf{18.} 在日照固定而出射功率减小的情况下，整个
气柱——包括地表在内——必须升温，直到
收支重新平衡。
\textbf{19.} 中心饱和，翼部活跃：气体每加倍一次，
不透明区域便以相近的增量展宽，因而
辐射强迫近似为对数形式——起作用的是第 11 题中的
翼部。
\textbf{20.} 不对称伸缩频率的平方正比于 $(1/m_{
\text{O}} + 2/m_{\text{C}})$：$^{13}$C 与 $^{12}$C 之比
为 $\sqrt{0.2163/0.2292} = 0.971$——下降 $2.9\%$，与
$2283/2349$ 相符。
\textbf{21.} 量子化使每种同位素分子各自拥有一组
尖锐的梳状谱线；把激光停驻于其中一组，只能数出一种同位素
——若吸收是经典连续谱，这是不可能的。
\textbf{22.} 随着 CO$_2$ 增加，其 $^{13}$C 比例\emph{下降}
（即苏斯效应）：新增的碳在同位素组成上偏轻——
来自植物、长期埋藏，即化石碳。大气携带了
其来源的特征。
\textbf{23.} 金星：同样的 \qty{15}{\micro m} 量子，作用于
跨越 $200000$ 倍的气柱，维持比烤箱更热的地表
——该机制没有内在的上限。
\textbf{24.} 火星几乎全是 CO$_2$，却仍然寒冷，表明
\emph{气柱总量与压强}（它们使谱线展宽），而非
气体比例，决定其强弱。
\textbf{25.} 一个 \qty{83}{meV} 的量子，波长为 \qty{15}{\micro m}，
在 \qty{288}{K} 行星的热肩部被吸收，热化
于纳秒内完成，从寒冷的高空再辐射：谐振子
阶梯，与维恩位移定律相权衡，正是这套机制
使万分之四的分子主宰着气候。
\end{solution}
